\section*{第 \ref{ch:b3:forms} 章 --- 微分形式与
斯托克斯定理}

\begin{solution}{exo:b3:forms:1}
展开并消去重复的因子：
\[
\omega\wedge\eta = (x\,\dd y - z\,\dd x)\wedge(\dd x +
y\,\dd z)
= -x\,\dd x\wedge\dd y + xy\,\dd y\wedge\dd z +
yz\,\dd z\wedge\dd x
\]
（利用 $\dd y\wedge\dd x = -\dd x\wedge\dd y$ 与 $-zy\,\dd
x\wedge\dd z = yz\,\dd z\wedge\dd x$）。接下来 $\dd\omega =
\dd x\wedge\dd y - \dd z\wedge\dd x = \dd x\wedge\dd y +
\dd x\wedge\dd z$ 与 $\dd\eta = \dd y\wedge\dd z$。最后
\[
\dd(\omega\wedge\eta) = y\,\dd x\wedge\dd y\wedge\dd z +
z\,\dd y\wedge\dd z\wedge\dd x = (y + z)\,\dd x\wedge\dd
y\wedge\dd z
\]
（$\omega\wedge\eta$ 的第一项贡献 $\dd(-x)
\wedge\dd x\wedge\dd y = 0$；三个因子的轮换
都是偶置换）。莱布尼茨检验：$\dd\omega\wedge\eta =
(\dd x\wedge\dd y + \dd x\wedge\dd z)\wedge(\dd x + y\,\dd
z) = y\,\dd x\wedge\dd y\wedge\dd z$，而
$(-1)^1\omega\wedge\dd\eta = -(x\,\dd y - z\,\dd
x)\wedge\dd y\wedge\dd z = z\,\dd x\wedge\dd y\wedge\dd z$；
两者之和相符。
\end{solution}

\begin{solution}{exo:b3:forms:2}
$\dd f = \sum_i\partial_if\,\dd x_i$ 的系数即
$\nabla f$ 的各分量。对于功形式，
\[
\dd\omega_F = (\partial_yF_3 - \partial_zF_2)\,\dd y\wedge\dd
z + (\partial_zF_1 - \partial_xF_3)\,\dd z\wedge\dd x +
(\partial_xF_2 - \partial_yF_1)\,\dd x\wedge\dd y =
\sigma_{\operatorname{curl}F},
\]
即旋度的通量形式（合并
$\sum_i\dd F_i\wedge\dd x_i$ 的六项）。对于通量形式，
$\dd\sigma_F = (\partial_xF_1 + \partial_yF_2 +
\partial_zF_3)\,\dd x\wedge\dd y\wedge\dd z$，即
散度。于是 $\dd^2f = 0$ 化为
$\sigma_{\operatorname{curl}\nabla f} = 0$，亦即
$\operatorname{curl}\nabla f = 0$，而 $\dd^2\omega_F = 0$
化为 $(\operatorname{div}\operatorname{curl}F)\,\dd
x\wedge\dd y\wedge\dd z = 0$：两个向量恒等式实为
同一恒等式 $\dd^2 = 0$，在两个次数上的体现。
\end{solution}

\begin{solution}{exo:b3:forms:3}
(a) 闭性即混合偏导数的对称性：
$\partial_y(2xy + z^2) = 2x = \partial_x(x^2)$,
$\partial_z(2xy + z^2) = 2z = \partial_x(2xz)$,
$\partial_z(x^2) = 0 = \partial_y(2xz)$。定义域 $\R^3$
是星形区域，故恰当（见\cref{thm:b3:forms:poincare}），
原函数为 $f = x^2y + xz^2$（验证 $\dd f$）。
(b) $\dfrac{x\,\dd x + y\,\dd y}{x^2 + y^2} =
\tfrac12\,\dd\log(x^2 + y^2)$：在整个
$\R^2\setminus\{0\}$ 上恰当（从而闭） --- 与角形式相近的径向
形式并无妨碍。
(c) 在 $\{x > 0\}$ 上，该形式即 $\omega_\theta$，是闭的
（见\cref{exo:b3:forms:4}）；半平面是凸的，故
在其上恰当，且确有 $f = \arctan(y/x)$ 满足 $\dd f
= \frac{-y\,\dd x + x\,\dd y}{x^2 + y^2}$。在
半平面上恰当，在去心平面上不恰当：
障碍在于洞，而不在公式本身。
\end{solution}

\begin{solution}{exo:b3:forms:4}
闭性：记 $\rho^2 = x^2 + y^2$，
\[
\partial_x\Bigl(\frac{x}{\rho^2}\Bigr) = \frac{\rho^2 -
2x^2}{\rho^4} = \frac{y^2 - x^2}{\rho^4} =
\partial_y\Bigl(\frac{-y}{\rho^2}\Bigr),
\]
于是 $\dd\omega_\theta = \bigl(\partial_x(x/\rho^2) -
\partial_y(-y/\rho^2)\bigr)\,\dd x\wedge\dd y = 0$。取
$\gamma(t) = (r\cos2\pi t, r\sin2\pi t)$：
\[
\gamma^*\omega_\theta = \frac{2\pi r^2(\cos^22\pi t +
\sin^22\pi t)}{r^2}\,\dd t = 2\pi\,\dd t,
\qquad\text{于是}\qquad
\int_\gamma\omega_\theta = 2\pi .
\] 若 $\omega_\theta$
恰当，则该积分应为零
（见\cref{prop:b3:forms:exactloop}）：它不是恰当的。若
$\R^2\setminus\{0\}$ 关于某点 $p$ 是星形的，
则由庞加莱引理（平移到 $p$ 处）可知每个
闭形式都恰当 --- 矛盾。直接地，对任意 $p
\neq 0$，连接 $p$ 与区域中的点 $-p$ 的线段
必经过 $0$：星形性在每一点都
不成立。
\end{solution}

\begin{solution}{exo:b3:forms:5}
对向量 $v_1, \dots, v_{n-1}$，将行列式
按其第一列展开，得
\[
\det(\nu, v_1, \dots, v_{n-1}) =
\sum_{i=1}^n(-1)^{i-1}\nu_i\,
\det\bigl(\text{行号} \neq i \text{ 的 } (v_1, \dots,
v_{n-1})\bigr)
= \sigma_\nu(v_1, \dots, v_{n-1}),
\]
由于 $(\dd x_1\wedge\dots\wedge\widehat{\dd
x_i}\wedge\dots\wedge\dd x_n)(v_1, \dots, v_{n-1})$ 恰为
第 $i$ 个余子式
（见\cref{def:b3:forms:wedge}）。将此应用于 $v_j =
\partial_j\gamma$：$(\gamma^*\sigma_\nu)(e_1, \dots,
e_{n-1}) = \det A$，其中 $A = (\nu\ \partial_1\gamma\ \cdots\
\partial_{n-1}\gamma)$。现 ${}^t\!AA$ 为分块对角阵：
$\langle\nu, \nu\rangle = 1$ 且 $\langle\nu,
\partial_j\gamma\rangle = 0$（$\nu$ 为法向，
$\partial_j\gamma$ 为切向），故 $(\det A)^2 =
\det({}^t\!AA) = \det G$；且 $\det A > 0$，对直接
参数化（这正是 $\nu$ 定向的含义）：
$\gamma^*\sigma_\nu = \sqrt{\det G}\,\dd
u_1\wedge\dots\wedge\dd u_{n-1}$，即格拉姆面积元。
对 $S^2$，$\nu(x) = x$，球面参数化为
$\gamma(\theta, \varphi) = (\sin\varphi\cos\theta,
\sin\varphi\sin\theta, \cos\varphi)$，定义在
$(0,2\pi)\times(0,\pi)$（直接的；它漏掉一条子午线，这是
一个零面积集）：$\det G = \sin^2\varphi$，于是
$\int_{S^2}\sigma_\nu =
\int_0^{2\pi}\!\!\int_0^\pi\sin\varphi\,\dd\varphi\,
\dd\theta = 4\pi$。
\end{solution}

\begin{solution}{exo:b3:forms:6}
给定的 $\omega$ 即为 $\sigma_\nu$，其中 $\nu(x) = x$，
定义在 $S^2$ 上，故 (a) 即刚才的计算：
$\int_{S^2}\omega = 4\pi$。(b) $\dd\omega = 3\,\dd
x\wedge\dd y\wedge\dd z$，由单位球上的斯托克斯公式得
$\int_{S^2}\omega = 3\operatorname{vol}(B^3)$，于是
$\operatorname{vol}(B^3) = \frac{4\pi}3$。一般地，
形式 $\sigma = \sum_i(-1)^{i-1}x_i\,\dd
x_1\wedge\dots\wedge\widehat{\dd x_i}\wedge\dots\wedge\dd
x_n$ 在 $S^{n-1}$ 上的限制为面积形式（取 $\nu = x$，
见\cref{exo:b3:forms:5}），$\dd\sigma = n\,\dd
x_1\wedge\dots\wedge\dd x_n$，由斯托克斯公式得
$\operatorname{area}(S^{n-1}) = n\operatorname{vol}(B^n)$，
--- 与伽马函数公式一致，
见\cref{thm:b3:product:ballvolume}。
\end{solution}

\begin{solution}{exo:b3:forms:7}
(a) 对通量形式 $\sigma_F$ 在
紧区域 $\Omega$ 上应用斯托克斯公式：
$\int_\Omega\operatorname{div}F\,\dd x =
\int_{\partial\Omega}\sigma_F$（内侧情形见\cref{exo:b3:forms:2}）。
确定边界上的被积式：对
$v_1, v_2$ 为 $\partial\Omega$ 上一点处的切向量，
由\cref{exo:b3:forms:5}的第一列展开式得
$\sigma_F(v_1, v_2) = \det(F, v_1, v_2)$；记 $F =
\langle F, \nu\rangle\nu + T$，其中 $T$ 为切向量；列
$T$ 落在 $v_1, v_2$ 张成的平面内，故 $\det(T,
v_1, v_2) = 0$，于是
$\sigma_F\vert_{\partial\Omega} = \langle F,
\nu\rangle\,\sigma_\nu = \langle F, \nu\rangle\,\dd S$：这就是
散度定理，其中 $\nu$ 为外法向
（外法向排在首位恰好给出诱导定向）。
(b) 对 $\omega_F$ 在
带边曲面 $S$ 上应用斯托克斯公式：$\dd\omega_F =
\sigma_{\operatorname{curl}F}$ 限制为
$\langle\operatorname{curl}F, \nu\rangle\,\dd S$，由同样的
对应关系；而在边界曲线上
$\gamma^*\omega_F = \langle F\circ\gamma,
\gamma'\rangle\,\dd t$，即 $\oint\langle F,
\tau\rangle\,\dd\ell$。上半球面，外
（径向）法向 $\nu$：在赤道点 $p = (1,0,0)$ 处，
曲面内向外的切向量为 $-e_3$；将其补成
正标架后可知赤道沿
从上方看为逆时针方向（$+e_3$），这即右手
法则，恒等式两侧约定一致。
\end{solution}

\begin{solution}{exo:b3:forms:8}
应用散度定理（\cref{exo:b3:forms:7}(a)，
其证明与维数无关），取 $F = u\nabla v$：
$\operatorname{div}(u\nabla v) = u\,\Delta v + \langle\nabla
u, \nabla v\rangle$，且 $\langle F, \nu\rangle =
u\,\partial_\nu v$：即第一恒等式。交换 $u, v$ 并
相减，对称项相消，即第二恒等式。若
$\Delta u = 0$ 于 $\Omega$ 上，且 $u = 0$ 于 $\partial\Omega$ 上：
由第一恒等式，取 $v = u$ 得
$\int_\Omega\norm{\nabla u}^2 = 0$，于是 $\nabla u \equiv 0$，
$u$ 在每个连通分支上为常数；每个分支的
闭包与 $\partial\Omega$ 相交（由有界性），而此处 $u = 0$：
$u \equiv 0$。两个调和函数若边界值相同，
则二者之差即为此类 $u$，故二者重合 --- 这是
狄利克雷问题解的唯一性，与
在\cref{ch:b3:conformal}中圆盘上的存在性理论互补。
\end{solution}

\begin{solution}{exo:b3:forms:9}
由\cref{thm:b3:forms:pullback}(c)，
$\varphi^*(f\,\dd y_1\wedge\dots\wedge\dd y_n) =
(f\circ\varphi)\,\det(D\varphi)\,\dd
x_1\wedge\dots\wedge\dd x_n$，故拉回恒等式为
\[
\int_U(f\circ\varphi)\,\det(D\varphi)\,\dd x = \int_Vf\,\dd
y .
\]
由于 $\det D\varphi > 0$ 处处成立，$\det D\varphi =
\abs{\det D\varphi}$，这正是
换元公式
（\cref{thm:b3:product:changeofvar}）对连续、
紧支撑的被积函数：两条陈述彼此等价，
绝对值则被移入\emph{假设}：
定向。对定向反转的 $\varphi$，形式
恒等式带上整体负号（形式感知
定向），而测度公式保留 $\abs{\det}$
（测度不感知）：同一雅可比行列式的两种记法。
\end{solution}

\begin{solution}{exo:b3:forms:10}
壳层 $A = \{r_1 \leq \norm x \leq r_2\}$ 是一个紧的
带边 $3$ 维子流形，边界为 $S_{r_2}\cup S_{r_1}$；
诱导定向在
$S_{r_2}$ 上为通常的球面定向（自 $A$ 向外即远离 $0$），而在
$S_{r_1}$ 上则为\emph{相反}的定向（自 $A$ 向外即指向
$0$）。由斯托克斯公式及 $\dd\omega = 0$：
\[
0 = \int_A\dd\omega = \int_{S_{r_2}}\omega -
\int_{S_{r_1}}\omega .
\]
立体角形式：取 $\rho = \norm x$ 及 $\sigma =
x\,\dd y\wedge\dd z + y\,\dd z\wedge\dd x + z\,\dd
x\wedge\dd y$，$\omega = \rho^{-3}\sigma$，且
\[
\dd\omega = \sum_i\partial_i\bigl(x_i\rho^{-3}\bigr)\,
\dd x\wedge\dd y\wedge\dd z
= \bigl(3\rho^{-3} - 3\rho^{-5}\cdot\rho^2\bigr)\dd
x\wedge\dd y\wedge\dd z = 0 .
\]
在 $S_r$ 上，对切向量 $v_1, v_2$：$\omega(v_1, v_2) =
r^{-3}\det(x, v_1, v_2) = r^{-3}\,r\det(\nu, v_1, v_2) =
r^{-2}\,\dd S(v_1, v_2)$，故 $\int_{S_r}\omega =
r^{-2}\cdot4\pi r^2 = 4\pi$：与 $r$ 无关，与形变
所预言的一致，且非零 --- 故 $\omega$ 为闭形式，但
在 $\R^3\setminus\{0\}$ 上不是恰当形式（恰当形式的积分须为
$0$，由斯托克斯公式，在无边界的 $S_r$ 上亦然）。这是高斯
定律：单位电荷之场穿过任一包围球面的通量
恒为 $4\pi$，与半径无关。
\end{solution}

\begin{solution}{exo:b3:forms:11}
写出 $\omega = c\,\omega_\theta + \dd f$
（见\cref{exo:b3:forms:12}），取 $c =
\frac1{2\pi}\int_{S^1}\omega$。于是
\[
\int_\gamma\omega = c\int_\gamma\omega_\theta +
\int_\gamma\dd f
= c\cdot2\pi\operatorname{Ind}_\gamma(0) + 0
= n\int_{S^1}\omega,
\]
由\cref{prop:b3:forms:exactloop}及
环绕数的定义。唯一整数 $n$ 控制
挖去原点的平面上的全部周期：闭 $1$-形式无法区分两
条环绕数相同的闭曲线。
\end{solution}

\begin{solution}{exo:b3:forms:12}
设 $\alpha = \omega - c\,\omega_\theta$，则它为闭形式，且
$\int_{S^1}\alpha = 0$，由 $c$ 的选取得出
（$\int_{S^1}\omega_\theta = 2\pi$）。用
极坐标映射 $\Phi(\rho, \theta) = (\rho\cos\theta, \rho\sin\theta)$ 作拉回，
是满射局部微分同胚 $(0, \infty)\times\R \to
U$：$\Phi^*\alpha$ 为闭形式
（由\cref{thm:b3:forms:pullback}(b)），它在\emph{凸}开集
$(0,\infty)\times\R$ 上，故为恰当形式
（由\cref{thm:b3:forms:poincare}）：$\Phi^*\alpha = \dd g$。
对固定的 $\rho$：$g(\rho, \theta + 2\pi) - g(\rho, \theta)
= \int_\theta^{\theta + 2\pi}\partial_\theta g\,\dd s$ 是
$\alpha$ 沿半径为
$\rho$ 的圆周的积分，它等于 $\int_{S^1}\alpha = 0$（
两圆之间的圆环是紧带边曲面；
仿\cref{exo:b3:forms:10}，用斯托克斯公式降一维）。
于是 $g$ 为 $2\pi$ 周期的（关于 $\theta$），从而
诱导出良定函数 $f$（定义在 $U$ 上），满足 $f\circ\Phi = g$；
$f$ 光滑（因 $\Phi$ 为局部微分同胚），且
$\Phi^*(\dd f) = \dd g = \Phi^*\alpha$ 推出 $\dd f =
\alpha$。于是 $\omega = c\,\omega_\theta + \dd f$。
唯一性：沿 $S^1$ 积分可确定 $c$，因
恰当形式的周期为零；$f$ 则至多相差一个加法
常数。映射 $[\omega] \mapsto
\frac1{2\pi}\int_{S^1}\omega$ 因而是线性
同构 $H^1(\R^2\setminus\{0\}) \to \R$，
且 $\omega_\theta$ 的上同调类生成：洞
在上同调意义下恰为一维。
\end{solution}

\begin{solution}{pb:b3:forms:1}
\textbf{1.} $\dd\sigma = \sum_i(-1)^{i-1}\,\dd x_i\wedge\dd
x_1\wedge\dots\wedge\widehat{\dd x_i}\wedge\dots\wedge\dd
x_n$；将 $\dd x_i$ 移过前面的 $i-1$ 个因子
代价为 $(-1)^{i-1}$，恰好抵消前面的系数，因而各
$n$ 项均等于 $\dd x_1\wedge\dots\wedge\dd x_n$，故
$\dd\sigma = n\,\dd x_1\wedge\dots\wedge\dd x_n$。对
球体用斯托克斯公式：$\int_S\sigma = \int_{\bar B}\dd\sigma =
n\operatorname{vol}(\bar B) > 0$。

\textbf{2.} $n = 2$：$\sigma = x\,\dd y - y\,\dd x$；在
$\gamma(t) = (\cos t, \sin t)$ 上，$\gamma^*\sigma = (\cos^2t +
\sin^2t)\,\dd t = \dd t$，即弧长形式：
$\int_S\sigma = 2\pi = 2\operatorname{vol}(\bar B^2)$。
$n = 3$：$\sigma\vert_S$ 为面积形式
（由\cref{exo:b3:forms:5}，取 $\nu(x) = x$）：
$\int_S\sigma = 4\pi = 3\cdot\tfrac{4\pi}3$。两者均与
第 1 问相符。

\textbf{3.} 对 $\norm\varphi^2 = 1$ 求导：$2\langle
D\varphi(x)h, \varphi(x)\rangle = 0$，对每个 $h$ 成立，故
$\operatorname{im}D\varphi(x) \subseteq \varphi(x)^\perp$，为
超平面：$\operatorname{rk}D\varphi(x) \leq n - 1$。由于
$\dd\sigma$ 为微分 $n$-形式（第 1 问），逐点有
$\varphi^*(\dd\sigma)_x =
(D\varphi(x))^*(\dd\sigma)_{\varphi(x)} = 0$，由
\cref{prop:b3:forms:linearpullback}(2)：映入
球面的映射没有余地拉回体积形式。

\textbf{4.} 错误在于第二个“因此”：在
$(n-1)$ 维的 $S$ 上，每个 $(n-1)$-形式显然都
是闭形式（因不存在非零的 $n$-形式，
在 $(n-1)$ 维流形上），但 \emph{闭不等于恰当}，
而 $\int_S\dd\eta = 0$ 要求存在真正的原形式 $\eta$
定义在 $S$ 上。$\sigma$ 的限制恰恰
不是恰当形式 --- 其积分为 $n\operatorname{vol}(\bar B) \neq
0$ --- 这一非恰当性支撑了整个问题。

\textbf{5.} 对 $r^*\sigma$ 在球面上积分，
并以两种方式计算。由于 $r$ 逐点固定 $S$，
该积分等于 $\int_S\sigma = n\operatorname{vol}(\bar B)
\neq 0$。由于 $r$ 定义在球体上，斯托克斯公式
将同一积分化为 $\int_{\bar B}\dd(r^*\sigma) =
\int_{\bar B}r^*(\dd\sigma)$；又因 $r$ 取值于
球面，由第 3 问知该被积式为零。同一个
数，两个值：这样的收缩映射不存在。

\textbf{6.} 两个积分都通过直接
参数化 $\gamma$（$S$ 的各片段）计算
（见\cref{def:b3:forms:integral}）；由于 $r\circ\gamma =
\gamma$（参数化落在 $S$ 上，其中 $r$ 为
恒等映射），$\gamma^*(r^*\sigma) = (r\circ\gamma)^*\sigma =
\gamma^*\sigma$（由\cref{thm:b3:forms:pullback}(a)）：局部
被积式相同，任一单位分解都给出
$\int_Sr^*\sigma = \int_S\sigma$。

\textbf{7.} $r^*\sigma$ 是光滑的 $(n-1)$-形式，定义在
紧定向闭球 $\bar B$ 的某邻域上，其
在诱导定向下的边界为 $S$：由斯托克斯公式
（\cref{thm:b3:forms:stokes}）恰好给出
$\int_Sr^*\sigma = \int_{\bar B}\dd(r^*\sigma)$。

\textbf{8.} $\dd(r^*\sigma) = r^*(\dd\sigma)$
（\cref{thm:b3:forms:pullback}(b)），它由
第 3 问（取 $\varphi = r$）得其为零。将第
6--8 问串联起来：
\[
0 < n\operatorname{vol}(\bar B) = \int_S\sigma =
\int_Sr^*\sigma = \int_{\bar B}\dd(r^*\sigma) = 0 :
\]
荒谬。\emph{不存在从 $\bar B^n$
到 $S^{n-1}$ 上的光滑收缩}（$n \geq 2$）。

\textbf{9.} 设连续映射 $r\colon\intcc{-1}1\to\{-1,1\}$
满足 $r(\pm1) = \pm1$，则它会把连通集映到
不连通的 $\{-1, 1\}$ 上，这不可能：连通集的连续像
仍连通 --- 等价地，介值定理
将迫使 $r$ 取到值
$0$。同一陈述在任意维数下正是
第 II 部分的内容；连通性是 $1$ 维
情形下上同调障碍 $\int_S\sigma \neq 0$ 的投影。

\textbf{10.} $x \mapsto \norm{g(x) - x}$ 连续，且
处处 $> 0$，定义于紧集 $\bar B$ 上：其最小值
$\delta$ 可以取到，故 $> 0$。

\textbf{11.} $\norm{x + tu}^2 = 1$ 即 $t^2 + 2t\langle x,
u\rangle + \norm x^2 - 1 = 0$，其根为
\[
t_\pm = -\langle x, u\rangle \pm \sqrt{\langle x,
u\rangle^2 + 1 - \norm x^2} .
\]
两根之积为 $\norm x^2 - 1 \leq 0$：两根分居
$0$ 两侧（或其一为零），故射线参数 ---
即非负根 --- 为 $t(x) = t_+$。

\textbf{12.} 若被开方式在 $x \in \bar B$ 处为零：
两项均非负，故 $\norm x = 1$ 且
$\langle x, u(x)\rangle = 0$，即 $\langle x, x -
g(x)\rangle = 0$，即 $\langle x, g(x)\rangle = 1$。由
柯西–施瓦茨不等式，$1 = \langle x, g(x)\rangle \leq
\norm x\,\norm{g(x)} \leq 1$：处处取等，这迫使
$g(x)$ 与 $x$ 共线，范数为 $1$，且方向相同：
$g(x) = x$ --- 与假设矛盾。因此被开方式连续且
$> 0$，在 $\bar B$ 上，从而存在 $c > 0$
作为下界，在其上及某邻域内亦然（由一致连续性）。在该
邻域上，$u$ 光滑（$\norm{x - g(x)} \geq
\delta/2$，必要时缩小邻域），被开方式保持 $\geq
c/2$，且平方根在 $\intoo0\infty$ 上光滑，故 $r$
在 $\bar B$ 附近光滑。

\textbf{13.} $\norm{r(x)} = 1$，由 $t(x)$ 的构造：
$r(\bar B) \subseteq S$。对 $x \in S$：$\langle x,
u(x)\rangle = \frac{1 - \langle x, g(x)\rangle}{\norm{x -
g(x)}} \geq 0$（再次应用柯西–施瓦茨不等式），且 $\norm x =
1$ 时，被开方式化为 $\langle x, u\rangle^2$，其
平方根即 $\langle x, u\rangle$ 本身（它 $\geq
0$）：$t(x) = 0$，$r(x) = x$。因此 $r$ 是一个
从球到球面的光滑收缩 --- 与
第 II 部分矛盾。\emph{$\bar B$ 上每个光滑
自映射都有不动点。}

\textbf{14.} $\varepsilon > 0$ 的取法同第 10 题。
多项式构成 $\mathcal C(\bar B,
\R)$ 的一个子代数，包含常值函数且分离点（$x
\mapsto x_i$ 即可），故由斯通–魏尔斯特拉斯
（\cref{thm:b3:complete:stoneweierstrass}）逼近每个
坐标：取多项式 $p_i$，使 $\sup_{\bar
B}\abs{p_i - f_i} < \frac{\varepsilon}{2\sqrt n}$；则
向量映射 $p = (p_1, \dots, p_n)$ 满足
$\sup_{\bar B}\norm{p - f} \leq
\bigl(\sum_i\sup_{\bar B}\abs{p_i - f_i}^2\bigr)^{1/2} <
\varepsilon/2$。

\textbf{15.} 在 $\bar B$ 上：$\norm p \leq \norm f +
\frac\varepsilon2 \leq 1 + \frac\varepsilon2$，故
$\norm{g} = \frac{\norm p}{1 + \varepsilon/2} \leq 1$：
$g(\bar B) \subseteq \bar B$。又 $\norm{g - p} =
\frac{\varepsilon/2}{1 + \varepsilon/2}\norm p \leq
\frac\varepsilon2$，从而 $\norm{g - f} \leq \norm{g - p} +
\norm{p - f} < \varepsilon$ 在 $\bar B$ 上成立。

\textbf{16.} $g$ 为多项式，故光滑，且把
$\bar B$ 映入自身：由第 III 部分，存在 $x_0 = g(x_0)$。
于是 $\norm{f(x_0) - x_0} = \norm{f(x_0) - g(x_0)} <
\varepsilon = \min_{\bar B}\norm{f - \operatorname{id}}$：
矛盾。\emph{每个连续映射 $\bar B^n \to
\bar B^n$ 都有不动点。}

\textbf{17.} \emph{开球：}$f(x) = \frac{x + e_1}2$ 把
$B$ 映入 $B$（严格地 $\norm{f(x)} < 1$），其唯一不动点
$e_1$ 位于球面上：紧性不再成立。
\emph{球面：}对径映射 $x \mapsto -x$ 没有
不动点；$S$ 紧，但具有“错误”的
拓扑 --- 它恰是第 II 部分所说的不可收缩的情形。
\emph{环域：}绕任一角 $\not\equiv 0$ 的旋转
没有不动点；中间的洞使旋转得以“躲藏” --- 凸性（更
准确地说，类球的拓扑）丧失。布劳威尔定理
实际上针对 \emph{紧凸}集，
正如第 18 题所利用的。

\textbf{18.} 对 $x \in \Delta$：某个 $x_j > 0$，故 $(Ax)_i
\geq A_{ij}x_j > 0$，对每个 $i$；从而 $\norm{Ax}_1 > 0$，
于是 $T(x) = Ax/\norm{Ax}_1$ 有定义、连续，且
落在 $\Delta$ 中（各分量为正且和为 $1$）。
$\Delta$ 是凸紧集，在仿射超平面中具有非空内部，
即 $\{\sum x_i = 1\} \cong \R^{n-1}$；
自其重心出发的径向映射 --- 每条射线
与 $\partial\Delta$ 恰交于一点（由凸性与紧性），
且相应的规范函数连续 ---
是一个同胚 $\Delta \to
\bar B^{n-1}$。通过它搬运布劳威尔定理：$T$ 有
不动点 $x^*$，即 $Ax^* = \lambda x^*$，其中 $\lambda
= \norm{Ax^*}_1 > 0$；且 $x^* = Ax^*/\lambda$ 的各分量严格为正，
由开头的计算得出。\emph{正矩阵
有正特征向量。}

\textbf{19.} 由第 18 题，$A\pi = \lambda\pi$，$\pi \in
\Delta$，$\pi > 0$。对各分量求和：$\sum_i(A\pi)_i =
\sum_j\pi_j\sum_iA_{ij} = \sum_j\pi_j = 1$（各列之和为
$1$），而 $\sum_i\lambda\pi_i = \lambda$。故 $\lambda = 1$
且 $A\pi = \pi$：一个平稳概率向量 ---
即马尔可夫链的平衡态，PageRank 的数学
核心。

\textbf{20.} 在 $S$ 上，$\norm{F_t(x)}^2 = \norm x^2 +
2t\langle x, v(x)\rangle + t^2\norm{v(x)}^2 = 1 + 0 + t^2$，
由切向性及 $\norm v = 1$，$F_t(S) \subseteq
\sqrt{1+t^2}\,S$。

\textbf{21.} 固定一个由直接参数化
$\gamma$ 组成的有限图册，以及与 $t$ 无关的单位分解。
在一个坐标卡中，$F_t\circ\gamma = \gamma + t(v\circ\gamma)$，故
$(F_t\circ\gamma)^*\sigma$ 的每个系数都是若干
乘积之和，每个乘积含一个因子 $(x_i + tv_i)\circ\gamma$（关于
$t$ 为仿射）及一个 $(n-1)\times(n-1)$ 行列式，
其元素关于 $t$ 为仿射：它是关于 $t$ 的次数 $\leq n$ 的多项式，其
系数为坐标卡变量的光滑函数。乘以
与 $t$ 无关的单位分解函数并逐项积分：
$P(t) = \sum_{k=0}^n c_kt^k$，是一个多项式。

\textbf{22.} \emph{单射性：}$v$ 在 $S$ 上利普希茨连续（在紧集上光滑），
设常数为 $L$；当 $\abs t <
1/L$，$\norm{F_t(x) - F_t(y)} \geq (1 - \abs
tL)\norm{x - y} > 0$。\emph{微分同胚：}$G_t =
F_t/\sqrt{1 + t^2}$ 把 $S$ 映到 $S$；在坐标卡中，其雅可比矩阵
一致收敛于 $G_0 = \operatorname{id}$ 的对应矩阵，
$t \to 0$ 时，故对小的 $t$，它们可逆，且 $G_t$ 是
局部微分同胚（坐标卡中的\cref{thm:b3:submanifolds:ift}），
是单射，像集为开集（局部微分同胚），且像集
在\emph{连通}的 $S^{n-1}$（$n \geq 2$）中为紧集：
像集 $= S$，故 $G_t$ 是 $S$ 的微分同胚。
\emph{尺度变换：}在 $s_\lambda(x) = \lambda x$ 下，每个系数
$x_i$ 乘以 $\lambda$，$n - 1$ 个
微分中的每一个也乘以 $\lambda$：$s_\lambda^*\sigma =
\lambda^n\sigma$。由于 $F_t = s_{\sqrt{1+t^2}}\circ G_t$：
\[
P(t) = \int_SG_t^*\bigl(s_{\sqrt{1+t^2}}^{\,*}\sigma\bigr)
= (1 + t^2)^{n/2}\int_SG_t^*\sigma
= (1 + t^2)^{n/2}\int_S\sigma
\quad\text{ 对充分小的 } t,
\]
最后一个等号由\cref{lem:b3:forms:consistency} 得到（$G_t$
是 $S$ 的微分同胚，且对小的 $t$ 保定向，
其坐标卡中的雅可比行列式连续变化，
从不为零，且在 $t = 0$ 时为正）。

\textbf{23.} 若 $n$ 为奇数且存在光滑的单位切向场，
则由第 21--22 题，多项式
$P(t)/\int_S\sigma$（合理，因 $\int_S\sigma \neq 0$ 由第 1 题得出），
与 $(1 + t^2)^{n/2}$ 在 $0$ 附近一致；两个
在 $0$ 附近一致的光滑函数，其中一个为多项式，
则 $(1 + t^2)^{n/2}$ 在整个
$\R$ 上必等于该多项式。但若 $Q(t)^2 = (1 + t^2)^n$，且 $Q \in \R[t]$，
则由 $\R[t]$ 中的唯一因式分解
（\cref{ch:b3:rings}）给出不可约多项式 $t^2 + 1$ 在 $Q^2$ 中的重数为偶数，
而奇数重数 $n$ 出现于
$(1 + t^2)^n$ 中，矛盾。故
$S^{n-1}$ 上不存在光滑的单位切向场（$n$ 为奇数），即偶数维
球面。最后，若只有连续的处处非零切向场
$w$，则它也能产生一个光滑的单位切向场：将其延拓到一个邻域上，
取为 $\tilde w(x) = w(x/\norm x)$，逐分量磨光
（\cref{thm:b3:lp:regularization}）得到一个光滑的
$v_0$，使得 $\sup_S\norm{v_0 - \tilde w} <
\frac12\min_S\norm w$；再取切向投影 $v_1(x) = v_0(x)
- \langle v_0(x), x\rangle x$——在 $S$ 上，这至多改变 $v_0$
的法向分量，其本身至多为 $\norm{v_0 -
\tilde w}$（因 $\tilde w$ 为切向），故 $\norm{v_1 - \tilde
w} \leq 2\norm{v_0 - \tilde w} < \min\norm w$，且 $v_1$
在 $S$ 上处处非零；最后归一化：$v =
v_1/\norm{v_1}$ 即为光滑的单位切向场。因此，
在每个偶数维球面上，\emph{每个连续切向场
都有零点}：地球上每一阵风都会留下
一个平静点。

\textbf{24.} $\langle v(x), x\rangle = \sum_j(-y_jx_j +
x_jy_j) = 0$：切向；$\norm{v(x)} = \norm x = 1$：单位。
光滑性显然（线性映射）。故每个奇数维
球面都带有一个光滑的单位切向场——沿复直线乘以
$\iu$——而第 23 题的
障碍恰好取决于维数的奇偶性。在
多项式论证中，对偶数 $n$，函数 $(1 +
t^2)^{n/2}$ \emph{确实是}多项式，并不产生矛盾：证明
并非仅仅不适用，
其结论确实为假，正如 $v$ 所见证的。

\textbf{25.} 设 $f$ 在 $\bar B$ 上处处不为零。则
$g(x) = -f(x)/\norm{f(x)}$ 是连续映射 $\bar B \to S
\subseteq \bar B$，由布劳威尔不动点定理（第 16 题）得 $x^* =
g(x^*)$。由于 $g$ 取值于 $S$，故 $x^* \in S$，且
\[
\langle f(x^*), x^*\rangle = \bigl\langle f(x^*),\
-\tfrac{f(x^*)}{\norm{f(x^*)}}\bigr\rangle
= -\norm{f(x^*)} < 0,
\]
与边界假设矛盾。故 $f$ 有零点。
满射性：任给 $y \in \R^n$，对 $f(x) =
F(x) - y$ 在球 $\bar B(0, R)$ 上应用上述结论，其中 $R$ 足够大，使得
$\langle F(x), x\rangle \geq \norm y\,\norm x$ 在半径为
$R$ 的球面上（强制性）；则 $\langle f(x), x\rangle =
\langle F(x), x\rangle - \langle y, x\rangle \geq 0$ 在其上
（由柯西–施瓦茨不等式），于是重新标度后的命题给出
$f$ 的零点：$F(x) = y$。每个强制连续场都是满射
——变分存在性定理在不借助拓扑度时的拓扑影子，
由纯拓扑直接给出。
\end{solution}
